\documentclass[11pt,a4paper]{amsart}
\usepackage[T1]{fontenc}
\usepackage{amsmath,amssymb,amsthm,mathtools,enumitem,microtype}
\usepackage[margin=2.7cm]{geometry}
\usepackage{xcolor}
\definecolor{linkcol}{rgb}{0.10,0.20,0.50}
\usepackage[colorlinks=true,linkcolor=linkcol,citecolor=linkcol,urlcolor=linkcol]{hyperref}
\theoremstyle{plain}
\newtheorem{theorem}{Theorem}[section]
\newtheorem{proposition}[theorem]{Proposition}
\newtheorem{lemma}[theorem]{Lemma}
\newtheorem{corollary}[theorem]{Corollary}
\theoremstyle{definition}
\newtheorem{definition}[theorem]{Definition}
\newtheorem{example}[theorem]{Example}
\newtheorem{remark}[theorem]{Remark}
\newcommand{\RR}{\mathbb R}
\newcommand{\QQ}{\mathbb Q}
\newcommand{\CC}{\mathbb C}
\newcommand{\NN}{\mathbb N}
\newcommand{\one}{\mathbf 1}
\newcommand{\Irr}{\operatorname{Irr}}
\newcommand{\Hom}{\operatorname{Hom}}
\newcommand{\im}{\operatorname{im}}
\newcommand{\indm}{\operatorname{ind}_{-}}
\newcommand{\Rad}{\operatorname{rad}}
\newcommand{\Sh}{\operatorname{Sh}}
\newcommand{\Vect}{\operatorname{Vect}}
\newcommand{\Hess}{\operatorname{Hess}}
\newcommand{\Def}{\operatorname{Def}}
\newcommand{\Pcal}{\mathcal P}
\newcommand{\Mcal}{\mathcal M}
\begin{document}
\title[The Lorentzian Stokes coarea law]{The Lorentzian Stokes coarea law and its arithmetic reflection theorem}
\author{M.~Haufschild}
\address{Independent researcher}
\email{haufschildmichael@gmail.com}
\date{4 August 2026}
\subjclass[2020]{Primary 34M40, 32S60; Secondary 11M26, 14D24, 15A18, 18F20, 16G20}
\keywords{Stokes structure, irregular Riemann--Hilbert correspondence, microsupport,
vanishing cycle, Rees construction, wild character variety, central local system,
Riemann Xi function, Mertens theorem, rational adeles, Witt index}

\begin{abstract}
Let $B$ be the core adjacency matrix attached by Boalch--Yamakawa and Dou\c{c}ot to a finite
ramified Stokes spectrum, and let $C=2I-B$.  We prove that its full quadratic form is a positive
coarea energy minus one square of physical rank.  More precisely, finite Puiseux coefficients
generate canonical threshold blocks $A_e$ and
\[
 x^TCx=\sum_i x_i^2+\sum_e\ell_e\Bigl(\sum_i|A_e\cap i|x_i\Bigr)^2
              -(H-1)\Bigl(\sum_i r_ix_i\Bigr)^2.
\]
Consequently every coefficient-realized nonabelian Hodge core, and every type-$A$ marking of it,
is finite, affine, or Lorentzian: negative index plus nullity is at most one.  For an honest graph
this is equivalent to the spectral obstruction $\lambda_2(B)\leq2$, with simple affine boundary.

The same factorization survives physical-rank fusion over an arbitrary marked curve.  All local
channels remain positive and topology changes only the single common-rank square, so the joint
negative-index--nullity bound remains one for every genus and every nonnegative multi-Rees
specialization.  Capacity completion gives the exact finite/affine/Lorentz trichotomy, while
corona reduction is its Schur-complement recursion.

We also identify the finite wall labels analytically.  For every actual leading coefficient
$c\ne0$ and pole numerator $k>0$, the zero directions of
$\operatorname{Re}(ce^{-ik\theta})$ are exactly $2k$ simple walls.  Their explicit labels preserve
cyclic order by rotation through $\pi/k$ and record transverse orientation.  The induced
constructible boundary germ has one directed microstalk over each wall; normalized directed
microstalk mass equals mixed Hom irregularity.  Thus published Cartan polarization, analytic
Stokes walls, microlocal Morse groups, Rees channels, and corona elimination are shadows of one
Lorentzian coarea law.  We finally classify scalar twists invisible to every mixed Stokes channel:
for a coefficient-generated spectrum there is only one, the common center, and its determinant
weight is exactly the physical-rank functional in the coarea law.  The positive half-density
character $\delta_n(s)=\exp((2\operatorname{Re}s-1)\log n)$ therefore gives a pre-divisor
arithmetic central local system.  If $M(x)$ is positive, its determinant holonomy is
$\delta_n(s)^{M(x)}=\delta_{n^{M(x)}}(s)$.  At a zero of the completed Riemann function, holonomy
minus identity on the first nonzero jet is the rational-adelic cyclic transgression
\[
                 (\delta_n(\rho)^{M(x)}-1)c_\rho.
\]
Pairing its normalized value with the value at the reflected divisor
$\rho^\vee=1-\overline\rho$ produces the positive prime energy
\[
 \mathcal E_{N,M}(\rho)=
 \sum_{p\leq N}\frac{4\sinh^2(M(\Re\rho-\tfrac12)\log p)}p.
\]
For every $N\geq2$ and $M>0$, this energy vanishes exactly on the critical line.  More sharply,
if $M(\Re s_N-\tfrac12)\log N\to c$, then
\[
 \mathcal E_{N,M}(s_N)\longrightarrow
 4\int_0^1\frac{\sinh^2(cu)}u\,du,
\]
with vanishing below and divergence above this logarithmic window.  Thus the independently forced
Stokes rank is also the exact arithmetic renormalization exponent.  This realizes and quantifies
the obstruction in the unique rank port but does not assert its global vanishing on all $\Xi$
divisors.
\end{abstract}
\maketitle

\section{Introduction}

Ramified Stokes geometry contains two gluing mechanisms.  Locally, exponential factors split into
branches on a finite cover and are regrouped into Stokes circles by analytic continuation.
Globally, local Stokes spaces are fused after physical ranks have been identified.  Main result is
that both mechanisms obey one quadratic conservation law: all centered deformation channels form
a positive coarea energy, while every nonpositive contribution factors through physical rank.

This law applies directly to the established nonabelian Hodge diagram, not to a newly named
contact matrix.  For Stokes circles $i,j$, let $r_i$ be ramification degree and
$A_{ij}=\Irr\Hom(i,j)$.  Boalch--Yamakawa's core formula, in Dou\c{c}ot's normalization, is
\begin{equation}\label{eq:published-B}
 B_{ij}=A_{ij}-r_ir_j+\delta_{ij},\qquad C=2I-B,
\end{equation}
where the diagonal formula includes the correction $+1$
\cite{BoalchYamakawa2020,Doucot2021}.  We identify the quadratic form of this published matrix
$C$ with coefficient coarea and determine its complete inertia.  This identification is the
geometric content; the ensuing rank-one linear algebra is elementary.

Scalar ancestry also matters.  Hiroe proves, under multiplicity-free, regularity, and spectral
curve hypotheses, that an irregular rigidity index equals an Euler characteristic
\cite{Hiroe2019}; related local formulas involve Milnor numbers and intersection multiplicities
\cite{Hiroe2017}.  We do not present that scalar identity as new.  Our invariant is the entire
polarization on independent orbit charges.  Its coefficients are mixed $\Irr\Hom$ terms, and its
conclusion is an inertia law rather than one numerical evaluation.

Nor do we claim discovery of fission trees or their ultrametric contact law.  Fission trees classify
admissible deformation data, and the ultrametric constraints on nonabelian Hodge graphs are known
\cite{Boalch2025,Doucot2025}.  New point is that coefficient coarea turns those data into a
positive Gram form with exactly one subtractive rank channel, and that the matrix so obtained is
precisely \eqref{eq:published-B}.

Let $w^N=z$ be a common local cover and let
\[
                         q_a(w)=\sum_{k=1}^{K}c_{a,k}w^{-k}
\]
range through a finite $\mu_N$-stable set $A$ of distinct exponential factors.  For $a\ne b$,
let $k(a,b)$ be the largest exponent at which the coefficients differ and put
$h(a,b)=k(a,b)/N$; set $h(a,a)=0$.  Highest difference satisfies
\[
                         h(a,b)\leq\max\{h(a,c),h(c,b)\}.
\]
Hence threshold partitions form a canonical dendrogram.  Choose $H$ above all contacts.  If
$A_e$ are its blocks, $\ell_e>0$ their persistence lengths, $I=A/\mu_N$, and $y_a=x_{[a]}$, then
\begin{equation}\label{eq:intro-law}
 \boxed{Q_A(x)=\sum_{i\in I}x_i^2+
        \sum_e\ell_e\left(\sum_{a\in A_e}y_a\right)^2
        -(H-1)\left(\sum_{a\in A}y_a\right)^2.}
\end{equation}
No tree, cut family, or contact matrix appears among the hypotheses.  Branch blocks need not be
unions of Galois orbits.  Their exact descended vector is
$u_{e,i}=|A_e\cap i|$; assuming orbit saturation would lose ramified examples.

Formula \eqref{eq:intro-law} has four exact finite-dimensional readings.  Its positive linear
forms are principal conormals of cluster hyperplanes in charge space.  Their normalized
coefficients are directed microlocal Morse groups of the closed dominance germ on the analytic
boundary stratification.  They are graded coordinates of a multi-parameter filtered quadratic
family.  Their elimination is a corona Schur complement whose residual is an explicit linear
quotient.  These assertions do not identify the family with a Rees module of holonomic
$\mathcal D$-modules or the corona quotient with arbitrary sheaf-theoretic vanishing cycles; their
precise scope is stated in the corresponding sections.

For a genus-$g$ curve with $m\geq1$ marked points, identify local physical ranks and denote the
common value by $M$.  Define
\begin{equation}\label{eq:intro-fusion}
                      Q_{X,D}=\sum_{a=1}^{m}Q_{A_a}+2(1-g-m)M^2.
\end{equation}
Substituting \eqref{eq:intro-law} gives a positive-definite form plus one multiple of $M^2$.
Consequently
\begin{equation}\label{eq:intro-inertia}
 Q_{X,D}>0\ \text{on }\ker M\setminus\{0\},\qquad
 \indm(Q_{X,D})+\dim\Rad(Q_{X,D})\leq1.
\end{equation}
Here $\Rad(Q)$ is the radical of the polar bilinear form of $Q$; its dimension is the nullity of
$Q$.  Thus \eqref{eq:intro-inertia} is one joint bound, not two independent estimates.  This is
the global Stokes--microlocal Witt law.  Capacity completion replaces the auxiliary rank
coefficient by the intrinsic constrained minimum at $M=1$ and yields an exact
finite/affine/Lorentz trichotomy.  Every handle changes only this coefficient, by $-2$.

Analytic irregular Riemann--Hilbert theory supplies the surrounding realization theory
\cite{Sabbah,DKRH}.  It is not needed to count the boundary walls: we derive them from the actual
phase and prove their cyclic incidence explicitly.  For broader realizations we state a precise
transport theorem: any exact realization equivalence identifying this incidence diagram preserves
the quotient microstalk and every term of the coarea identity.

There is also a canonical arithmetic meaning of the rank line.  A scalar twist on each Stokes
circle acts on a mixed Hom channel by the quotient of the two scalars.  Since distinct
coefficient-generated circles have positive mixed irregularity, invisibility on all local channels
forces one common central scalar.  Its action on the determinant line of a dimension vector $x$ is
the scalar raised to $M(x)$.  Thus $M$ is simultaneously the unique nonpositive metric port and the
unique determinant weight invisible to relative Stokes data.  Section~\ref{sec:arithmetic}
constructs the half-density character in that center and identifies its holonomy defect on the
leading completed-Riemann jet with an explicit trace-one rational-adelic cyclic residue.  The
construction is unconditional.  Reflection pairing then gives a positive prime energy whose
finite vanishing locus is exactly the critical line and whose rank-scaled critical window has a
universal Mertens profile.  Annihilation at every divisor point remains equivalent to the Riemann
hypothesis and is not claimed.

Sections~\ref{sec:coefficients}--\ref{sec:global} prove coefficient and global laws.
Sections~\ref{sec:cellular}--\ref{sec:transport} construct categorical microlocal realization.
Sections~\ref{sec:rees}--\ref{sec:corona} treat Rees, capacity, and corona.
Section~\ref{sec:arithmetic} constructs the arithmetic central port.  Section~\ref{sec:deformation}
derives reduction dimension.  Final sections give marking legs, spectral consequences, novelty
boundaries, and falsifiers.

\section{Finite Puiseux spectra and canonical cuts}\label{sec:coefficients}

\begin{definition}[Equivariant finite spectrum]\label{def:spectrum}
Fix $N,K\geq1$ and a primitive $N$th root $\zeta$.  An equivariant finite Puiseux spectrum is a
finite set $A$, a permutation $\sigma:A\to A$ with $\sigma^N=1$, and coefficients
$c_{a,k}\in\CC$ such that
\[
                  c_{\sigma a,k}=\zeta^{-k}c_{a,k}
                  \qquad(a\in A,\ 1\leq k\leq K).
\]
Coefficient vectors are required to be pairwise distinct.
\end{definition}

For $a\ne b$, set
\[
 k(a,b)=\max\{k:c_{a,k}\ne c_{b,k}\},\qquad h(a,b)=k(a,b)/N,
\]
and put $k(a,a)=h(a,a)=0$.

\begin{lemma}[Highest-difference law]\label{lem:three-point}
For all $a,b,c\in A$,
\begin{equation}\label{eq:three-point}
                         k(a,b)\leq\max\{k(a,c),k(c,b)\}.
\end{equation}
Thus the largest of $k(a,b),k(a,c),k(b,c)$ occurs at least twice.
\end{lemma}

\begin{proof}
Above the maximum on the right, coefficients of $a,c$ agree and coefficients of $c,b$ agree.
Hence coefficients of $a,b$ agree.  Permuting the three variables excludes a unique largest value.
\end{proof}

For $t>0$, define
\begin{equation}\label{eq:threshold}
              a\sim_t b\quad\Longleftrightarrow\quad a=b\ \text{ or }\ h(a,b)<t.
\end{equation}

\begin{theorem}[Canonical coefficient dendrogram]\label{thm:dendrogram}
Every $\sim_t$ is an equivalence relation.  The resulting partitions $\Pcal_t$ coarsen with $t$.
Their finitely many blocks, ordered by inclusion, form a canonical rooted dendrogram.  For
distinct $a,b$, their smallest common block is born at $h(a,b)$.
\end{theorem}

\begin{proof}
If $a\sim_t b$ and $b\sim_t c$, Lemma~\ref{lem:three-point} gives $h(a,c)<t$; this proves
transitivity.  Monotonicity is immediate.  Only levels $k/N$ with $0\leq k\leq K$ occur.  Join
each block to the least strictly larger block and suppress repeated unary vertices.  Definition
\eqref{eq:threshold} identifies the birth level of the first common block.
\end{proof}

Choose $H>\max h(a,b)$.  Define the finite step integral
\begin{equation}\label{eq:partition-energy}
 E_H(y)=\int_0^H\sum_{C\in\Pcal_t}\left(\sum_{a\in C}y_a\right)^2dt,
\end{equation}
meaning interval length times the constant integrand between successive contact values.

\begin{theorem}[Coefficient coarea identity]\label{thm:coarea}
For every $y:A\to\RR$,
\begin{equation}\label{eq:coarea}
 E_H(y)=H\left(\sum_ay_a\right)^2-\sum_{a,b}h(a,b)y_ay_b.
\end{equation}
If $e$ runs through persistent non-root blocks $A_e$ and $\ell_e>0$ is persistence length, then
\begin{equation}\label{eq:cut-energy}
                         E_H(y)=\sum_e\ell_e\left(\sum_{a\in A_e}y_a\right)^2.
\end{equation}
Moreover $E_H(y)>0$ for $y\ne0$.
\end{theorem}

\begin{proof}
After expansion, the coefficient of $y_ay_b$ in \eqref{eq:partition-energy} is the length of
thresholds for which $a\sim_t b$, namely $H-h(a,b)$.  This proves \eqref{eq:coarea}.  Grouping the
step integral by persistent blocks gives \eqref{eq:cut-energy}.  At sufficiently small positive
$t$ all blocks are singletons, so a nonzero coordinate contributes on an interval of positive
length.
\end{proof}

\section{Orbit descent and the local Witt law}\label{sec:orbit}

Let $I=A/\langle\sigma\rangle$, write $[a]$ for the orbit of $a$, and put
\begin{equation}\label{eq:orbit-data}
 r_i=|i|,\qquad
 s_{ij}=\sum_{a\in i}\sum_{b\in j}h(a,b),\qquad
 M(x)=\sum_ir_ix_i.
\end{equation}

\begin{lemma}[Galois integrality]\label{lem:integrality}
Every $s_{ij}$ is a nonnegative integer.
\end{lemma}

\begin{proof}
Decompose $i\times j$ into cyclic pair orbits.  If one has length $L$ and leading difference
degree $k$, its stabilizer contains $\sigma^L$.  Equivariance gives $\zeta^{-Lk}=1$, so $N$ divides
$Lk$.  The pair orbit contributes $Lk/N\in\NN$.  Sum over pair orbits.
\end{proof}

\begin{definition}[Local polarized Stokes index]\label{def:local-index}
\begin{equation}\label{eq:local-index}
             Q_A(x)=\sum_i x_i^2+M(x)^2-\sum_{i,j}s_{ij}x_ix_j.
\end{equation}
\end{definition}

\begin{proposition}[Published Cartan identification]\label{prop:published-Cartan}
Let $A_{ij}=\Irr\Hom(i,j)=s_{ij}$ and define the published core adjacency by
\begin{equation}\label{eq:core-adjacency}
                 B_{ij}=A_{ij}-r_ir_j+\delta_{ij},\qquad C=2I-B.
\end{equation}
Then, entrywise,
\begin{equation}\label{eq:core-Cartan-entry}
                    C_{ij}=\delta_{ij}+r_ir_j-s_{ij},
\end{equation}
and its quadratic form is exactly
\begin{equation}\label{eq:published-Cartan-identification}
                              x^TCx=Q_A(x).
\end{equation}
Thus $Q_A$ is the published Cartan polarization of the Stokes circles, namely the cyclic
continuation orbits of the Puiseux spectrum.
\end{proposition}

\begin{proof}
The equality $A_{ij}=s_{ij}$ is the definition of irregularity of the Hom Stokes circle: on the
common cover it is the normalized sum of the highest pole orders over all ordered branch pairs.
Substitution of \eqref{eq:core-adjacency} into $C=2I-B$ gives
\eqref{eq:core-Cartan-entry}.  Summing against $x_ix_j$ gives
$\sum_i x_i^2+(\sum_i r_ix_i)^2-\sum_{i,j}s_{ij}x_ix_j$, which is
Definition~\ref{def:local-index}.  Formula \eqref{eq:core-adjacency} is precisely the published
core convention \cite{BoalchYamakawa2020,Doucot2021}.
\end{proof}

For a block $A_e$, define exact descended incidence and charge
\begin{equation}\label{eq:descended-cut}
        u_{e,i}=|A_e\cap i|,\qquad
        \nu_e(x)=\sum_i u_{e,i}x_i=\sum_{a\in A_e}x_{[a]}.
\end{equation}

\begin{theorem}[Local Stokes--microlocal Witt law]\label{thm:local-witt}
For every orbit charge $x$,
\begin{equation}\label{eq:local-witt}
 \boxed{Q_A(x)=\sum_i x_i^2+\sum_e\ell_e\nu_e(x)^2-(H-1)M(x)^2.}
\end{equation}
Hence $Q_A$ is positive definite on $\ker M$ and
\begin{equation}\label{eq:local-joint-inertia}
                 \indm(Q_A)+\dim\Rad(Q_A)\leq1.
\end{equation}
\end{theorem}

\begin{proof}
Lift $x$ to $y_a=x_{[a]}$.  Then
\[
 \sum_{i,j}s_{ij}x_ix_j=\sum_{a,b}h(a,b)y_ay_b,
 \qquad \sum_ay_a=M(x).
\]
Insert Theorem~\ref{thm:coarea} into Definition~\ref{def:local-index}.  On $\ker M$, the final
square vanishes and $\sum_i x_i^2$ is strictly positive.  A two-dimensional negative subspace
would meet this codimension-one hyperplane nontrivially.  The radical meets $\ker M$ trivially,
so it has dimension at most one.  Finally, a negative vector and a nonzero radical vector cannot
coexist: a nonzero linear combination with rank zero would retain negative energy because the
radical is orthogonal to every vector, contradicting positivity on $\ker M$.
\end{proof}

\begin{corollary}[Core Lorentz law]\label{cor:core-lorentz}
The generalized Cartan form of every core generated by a finite equivariant Puiseux spectrum is
positive definite, positive semidefinite of corank one, or Lorentzian.  In particular it has no
negative direction at affine boundary and no second negative or null channel.
\end{corollary}

\begin{proof}
Combine Proposition~\ref{prop:published-Cartan} with
\eqref{eq:local-joint-inertia}.
\end{proof}

\begin{remark}[Ramified descent]
Blocks $A_e$ may split a Galois orbit.  Formula \eqref{eq:descended-cut}, rather than an
orbit-saturation assumption, is required for ramified data.  When all blocks are orbit-saturated,
the usual weighted rooted-core cut vectors are recovered.
\end{remark}

\section{The global fused Witt law}\label{sec:global}

Let $D=\{p_1,\ldots,p_m\}$, $m\geq1$.  At $p_a$ choose a spectrum $A_a$, orbit space $I_a$,
local rank $M_a$, and local form $Q_a$.  Define
\begin{equation}\label{eq:fibre-product}
 V_{X,D}=\{(x_a)\in\bigoplus_a\RR^{I_a}:M_1(x_1)=\cdots=M_m(x_m)\}.
\end{equation}
Write $M$ for common rank and, for genus $g$, define
\begin{equation}\label{eq:fused-form}
             Q_{X,D}(x)=\sum_{a=1}^{m}Q_a(x_a)+2(1-g-m)M^2.
\end{equation}

\begin{proposition}[Quasi-Hamiltonian origin of fusion]\label{prop:qh-fusion}
For $G=\mathrm{GL}_M$, let $h_a$ be formal-centralizer dimension, $s_a$ directed
Stokes-root dimension, and $k_a$ prescribed conjugacy-class dimension.  Put
$Q_a^{\rm core}=M^2+h_a-s_a$ and $Q_a^{\rm mark}=Q_a^{\rm core}-k_a$.  Then
\begin{align}\label{eq:qh-dimension}
 \dim\Mcal_{X,D}
 &=2+(2g-2+m)M^2+\sum_a(s_a+k_a-h_a),\\
 Q_a^{\rm mark}&=M^2+h_a-s_a-k_a,
\end{align}
and therefore
\begin{equation}\label{eq:qh-fusion-identity}
       2-\dim\Mcal_{X,D}=\sum_aQ_a^{\rm mark}+2(1-g-m)M^2.
\end{equation}
The published diagram formula identifies $Q_a^{\rm core}$ with the coefficient form $Q_a$ in
\eqref{eq:fused-form}.  Thus \eqref{eq:fused-form} is the unmarked core polarization.  For a
type-$A$ marking, $Q_a^{\rm mark}$ is its leg-completed Cartan form
\cite{BoalchYamakawa2020,Doucot2021}; Section~\ref{sec:legs} proves that completion preserves the
same one-channel bound.
\end{proposition}

\begin{proof}
The local fission/Stokes factor has dimension $M^2+h_a+s_a$.  Fusion over a genus-$g$ surface,
reduction by the formal centralizers, insertion of the conjugacy classes, and restoration of the
two-dimensional scalar center give the first line of \eqref{eq:qh-dimension}
\cite{Boalch2014,BoalchYamakawa2015}.  The remaining equalities are algebraic rearrangements.
Core and type-$A$ Cartan identifications are the published diagram dimension formula.
\end{proof}

Choose local root heights $H_a$ and set
\begin{align}\label{eq:global-positive}
 P_{X,D}(x)&=\sum_a\left(\sum_{i\in I_a}x_{a,i}^2+
                    \sum_e\ell_{a,e}\nu_{a,e}(x_a)^2\right),\\
 \Gamma_{X,D}&=2(1-g-m)-\sum_a(H_a-1).\notag
\end{align}

\begin{theorem}[Global coefficient-generated Witt law]\label{thm:global-witt}
On $V_{X,D}$,
\begin{equation}\label{eq:global-decomposition}
                  \boxed{Q_{X,D}=P_{X,D}+\Gamma_{X,D}M^2.}
\end{equation}
The form $P_{X,D}$ is positive definite, and
\begin{equation}\label{eq:global-inertia}
 Q_{X,D}>0\ \text{on }\ker M\setminus\{0\},\qquad
 \indm(Q_{X,D})+\dim\Rad(Q_{X,D})\leq1.
\end{equation}
\end{theorem}

\begin{proof}
Sum Theorem~\ref{thm:local-witt} and collect common-rank squares.  Terminal sums make
$P_{X,D}$ positive definite.  Any two-dimensional negative subspace meets $\ker M$ nontrivially.
The radical meets $\ker M$ trivially, so its dimension is at most one.  If a negative vector and a
nonzero radical vector coexisted, their span would contain a nonzero rank-zero vector; since the
radical is orthogonal to all vectors, that combination would have negative energy, contradiction.
\end{proof}

\begin{corollary}[Handle and sewing laws]\label{cor:handle}
Adding one handle subtracts $2M^2$ and changes no local channel.  Joining two marked curves away
from marked points adds their scalar rank coefficients and then subtracts $2$.
\end{corollary}

\begin{proof}
Use \eqref{eq:fused-form} and $\chi(X\#Y)=\chi(X)+\chi(Y)-2$.
\end{proof}

\section{Analytic-wall Stokes categories and vanishing cycles}\label{sec:cellular}

Let $\mathsf C$ be the walking cospan $L\to W\leftarrow R$ and put
\[
                         \Sh_{\rm germ}=\operatorname{Fun}(\mathsf C,\Vect_{\QQ}).
\]
This is an abelian category.  For $S=L$ or $R$, write
$\rho_S(F):F(S)\to F(W)$ and define
\begin{equation}\label{eq:micro-functor}
                         \mu_S(F)=F(W)/\im\rho_S(F).
\end{equation}

\begin{proposition}[Microlocal quotient]\label{prop:microfunctor}
Equation \eqref{eq:micro-functor} defines a functor
$\mu_S:\Sh_{\rm germ}\to\Vect_{\QQ}$.  For every $F$,
\begin{equation}\label{eq:nearby-vanishing}
       F(S)\xrightarrow{\rho_S(F)}F(W)\longrightarrow\mu_S(F)\longrightarrow0
\end{equation}
is exact.
\end{proposition}

\begin{proof}
For a natural transformation $\alpha:F\to G$, naturality gives
$\alpha_W(\im\rho_S(F))\subseteq\im\rho_S(G)$.  Hence $\alpha_W$ induces a unique quotient map.
Identity and composition descend.  Exactness and surjectivity are the defining properties of a
quotient vector space.
\end{proof}

For a finite-dimensional vector space $V$ and $\varepsilon\in\{0,1\}$, define the exponential
germ $E(V,\varepsilon)$ by
\[
 E(W)=V,\qquad (E(L),E(R))=
 \begin{cases}(V,0),&\varepsilon=1,\\(0,V),&\varepsilon=0,
 \end{cases}
\]
with identity and zero incidence maps.

\begin{corollary}[One oriented vanishing cycle]\label{cor:one-oriented}
If $V\ne0$, exactly one of $\mu_L(E(V,\varepsilon))$ and
$\mu_R(E(V,\varepsilon))$ is zero; the other is canonically $V$.
\end{corollary}

\begin{proof}
The cokernel of $V\xrightarrow{1}V$ is zero; the cokernel of $0\to V$ is $V$.
\end{proof}

Fix $c\in\CC^\times$ and $k>0$.  On the oriented boundary circle of the common $w$-cover,
$\RR/2\pi\mathbb Z$, use the dominance-wall convention
\begin{equation}\label{eq:analytic-phase}
 f_{c,k}(\theta)=\operatorname{Re}(ce^{-ik\theta})
  =|c|\cos(\arg c-k\theta).
\end{equation}
These are the analytic walls separating the two signs of exponential dominance.

\begin{proposition}[Complete analytic wall classification]\label{prop:analytic-walls}
Put $\tau_+=\pi/2$ and $\tau_-=-\pi/2$.  The zero set of \eqref{eq:analytic-phase} is
\begin{equation}\label{eq:analytic-wall-labels}
 \mathsf W_k(c)=\left\{
 \theta_{j,\varepsilon}=
 \frac{\arg c-\tau_\varepsilon}{k}+\frac{2\pi j}{k}\quad(\bmod\,2\pi):
 0\leq j<k,\ \varepsilon\in\{+,-\}\right\}.
\end{equation}
All $2k$ classes in \eqref{eq:analytic-wall-labels} are distinct and exhaust the zero set.  Every
wall is simple, with
\[
 f'_{c,k}(\theta_{j,+})=k|c|,
 \qquad f'_{c,k}(\theta_{j,-})=-k|c|.
\]
Consequently $\{0,\ldots,k-1\}\times\{+,-\}$ is canonically bijective with the analytic Stokes
walls, and its sign coordinate is their transverse orientation.  If
\[
 (j,+)^+=(j,-),\qquad (j,-)^+=(j+1\bmod k,+),
\]
then
\begin{equation}\label{eq:analytic-successor}
                       \theta_{w^+}=\theta_w+\frac{\pi}{k}quad(\bmod,2\pi).
\end{equation}
Thus finite successor is the genuine positive cyclic successor, not an independently imposed
cellular order.
\end{proposition}

\begin{proof}
The equality $f_{c,k}(\theta)=0$ is equivalent on the angle circle to
$\arg c-k\theta=\pi/2$ or $-\pi/2$.  Multiplication by $k$ on
$\RR/2\pi\mathbb Z$ has, over each target, the $k$ distinct inverse images displayed in
\eqref{eq:analytic-wall-labels}.  The two fibers are disjoint because sine takes the values $1$
and $-1$ at their respective targets.  Finally
$f'_{c,k}(\theta)=k|c|\sin(\arg c-k\theta)$, which gives the asserted nonzero derivatives and
orientations.  Direct substitution in \eqref{eq:analytic-wall-labels} gives
\eqref{eq:analytic-successor}; at $j=k-1$ the added $2\pi$ vanishes on the angle circle.
\end{proof}

For two distinct coefficient branches $a,b$, put $k=k(a,b)$ and
$c=c_{a,k}-c_{b,k}$.  Highest-difference definition gives $k>0$ and $c\ne0$.  Proposition
\ref{prop:analytic-walls} therefore applies to every branch pair with no additional analytic-wall
hypothesis.  In particular, wall cardinality, cyclic incidence, and transverse orientation are
coefficient-derived.

Use Proposition~\ref{prop:analytic-walls} to label the actual analytic walls by
$\mathsf W_k=\{0,\ldots,k-1\}\times\{+,-\}$.  Place one arc between consecutive walls.  Its
entrance quiver has two arrows from each arc to its endpoints.  No two nonidentity arrows compose,
so the path category contains only identities and incidence arrows.  Define
\begin{equation}\label{eq:boundary-category}
     \Sh_{\rm St}(S^1;k)=
     \operatorname{Fun}(\operatorname{Path}(\mathsf W_k),\Vect_{\QQ}).
\end{equation}
The alternating exponential object has wall stalk $V$, alternating arc stalks $V,0$, and identity
or zero incidence maps.

\begin{definition}[Boundary microsupport and cycle]\label{def:cellular-cycle}
An oriented conormal $(w,S)$ lies in boundary microsupport when
$\mu_S(F|_w)\ne0$.  Its characteristic-cycle coefficient is
$\dim\mu_S(F|_w)$.  The directed cycle retains the $k$ conormals carrying ordered decay orientation.
\end{definition}

\begin{theorem}[Exponential microsupport]\label{thm:exponential-micro}
If $\dim V=d$, then the alternating degree-$k$ exponential object has exactly one microsupport
conormal over every wall, each nonzero cycle coefficient equals $d$, full cycle mass is $2kd$, and
directed decay mass is $kd$.
\end{theorem}

\begin{proof}
Each wall germ is $E(V,0)$ or $E(V,1)$, so the first two statements follow from
Corollary~\ref{cor:one-oriented}.  There are $2k$ walls and one directed wall in each of $k$
labelled antipodal pairs.
\end{proof}

\begin{corollary}[Analytic-wall characteristic cycle]\label{cor:analytic-cycle}
For two distinct coefficient branches $a,b$, put
$k=k(a,b)$ and $c=c_{a,k}-c_{b,k}$.  The preceding microsupport is supported on the actual
phase-zero set
\[
       \{\theta\in\RR/2\pi\mathbb Z:\Re(ce^{-ik\theta})=0\}.
\]
It has exactly one oriented conormal over every analytic wall, with coefficient $d$ after tensoring
by a rank-$d$ Hom space; its full and directed masses are respectively $2kd$ and $kd$.  The
orientation is the transverse derivative sign of the displayed phase.
\end{corollary}

\begin{proof}
Highest-difference definition gives $k>0$ and $c\ne0$.  Transport Theorem~\ref{thm:exponential-micro}
through the cyclic-order- and transverse-sign-preserving bijection of
Proposition~\ref{prop:analytic-walls}.
\end{proof}

Thus characteristic coefficients and vanishing ranks are dimensions of an actual quotient
functor, not names assigned to a numerical term.

\section{Microlocal Hom index and Galois descent}\label{sec:microlocal}

For branches $a,b$, let $E_{ab}$ be the exponential Hom object of degree $k(a,b)$.  If orbit $i$
carries a coefficient space $V_i$ of dimension $d_i$, tensor $E_{ab}$ by
$\Hom(V_{[a]},V_{[b]})$.

\begin{definition}[Normalized microlocal Hom mass]\label{def:micro-mass}
\begin{equation}\label{eq:micro-mass}
 \mathfrak m_{ij}(d)=\frac1N
 \sum_{a\in i}\sum_{b\in j}\sum_{\xi\in\operatorname{Decay}(a,b)}
       \dim\mu_\xi(E_{ab}\otimes\Hom(V_i,V_j)).
\end{equation}
\end{definition}

\begin{theorem}[Directed microlocal Hom index]\label{thm:micro-index}
For all orbit classes $i,j$,
\begin{equation}\label{eq:micro-index}
                         \mathfrak m_{ij}(d)=s_{ij}d_id_j.
\end{equation}
In particular, for one-dimensional coefficient spaces, normalized directed microstalk mass equals
mixed Hom irregularity.
\end{theorem}

\begin{proof}
Theorem~\ref{thm:exponential-micro} gives mass $k(a,b)d_id_j$ for one pair.  Sum and divide by
$N$ to obtain $s_{ij}d_id_j$.  Lemma~\ref{lem:integrality} proves integrality after normalization.
\end{proof}

Define $\langle x,x\rangle_\mu=\sum_{i,j}\mathfrak m_{ij}(1)x_ix_j$.

\begin{theorem}[Categorical microlocal coarea law]\label{thm:categorical-coarea}
\begin{align}\label{eq:categorical-coarea}
 Q_A(x)
 &=\sum_i x_i^2+M(x)^2-\langle x,x\rangle_\mu\notag\\
 &=\sum_i x_i^2+\sum_e\ell_e\nu_e(x)^2-(H-1)M(x)^2.
\end{align}
Hence analytic-wall microsupport coefficients, directed vanishing-cycle ranks, mixed Hom irregularities,
and threshold coarea coefficients are four presentations of the same polarization.
\end{theorem}

\begin{proof}
The first equality is Definition~\ref{def:local-index} and Theorem~\ref{thm:micro-index}; the second
is Theorem~\ref{thm:local-witt}.
\end{proof}

\section{Exact realization transport}\label{sec:transport}

Let $\mathcal A$ be an abelian category and
$\Phi:\mathcal A\overset\sim\longrightarrow\Sh_{\rm germ}$ an exact equivalence.  Set
$\Psi_W=\operatorname{ev}_W\Phi$ and $\phi_S=\mu_S\Phi$.

\begin{theorem}[Irregular-realization transport]\label{thm:transport}
For every $A\in\mathcal A$,
\begin{equation}\label{eq:transport-exact}
 \operatorname{ev}_S\Phi(A)\longrightarrow\Psi_W(A)
 \longrightarrow\phi_S(A)\longrightarrow0
\end{equation}
is exact and functorial.  If $\Phi(A)$ is exponential, precisely one transported oriented
vanishing object has full coefficient rank.  Naturally isomorphic realization equivalences give
isomorphic quotient microstalks and preserve every term of \eqref{eq:categorical-coarea}.
\end{theorem}

\begin{proof}
Apply Proposition~\ref{prop:microfunctor} to $\Phi(A)$.  Exactness and functoriality survive
composition.  The exponential statement is Corollary~\ref{cor:one-oriented}.  A natural
isomorphism of realization functors identifies stalk maps and therefore their cokernels.
\end{proof}

\begin{remark}[Analytic scope]
In Theorem~\ref{thm:transport}, an irregular-realization datum means exactly the displayed exact
equivalence with the analytic-wall germ category.  Relevant analytic constructions are developed
in \cite{Sabbah,DKRH,DKVanishing}; no existence or comparison theorem from them is used in the
proof of the coarea law.  Proposition~\ref{prop:analytic-walls} independently proves that the
incidence diagram is the stratification by the genuine leading exponential phase, not an abstract
polygon.  We do not infer a meromorphic connection or a realization equivalence from the
quadratic form.
\end{remark}

\section{Principal conormals and the multi-Rees law}\label{sec:rees}

For each block $A_e$, let $\mathcal H_e=\{x:\nu_e(x)=0\}$.  Its constant principal conormal is
$d\nu_e=u_e=(u_{e,i})_i$.  Put
$P_A(x)=\sum_i x_i^2+\sum_e\ell_e\nu_e(x)^2$.

\begin{proposition}[Principal-conormal Hessian]\label{prop:hessian}
\begin{equation}\label{eq:hessian}
                   \frac12\Hess(P_A)=I+\sum_e\ell_e\,d\nu_e\otimes d\nu_e.
\end{equation}
The sole subtractive Hessian term in $Q_A$ is $-(H-1)dM\otimes dM$.
\end{proposition}

\begin{proof}
Polarize the squares in Theorem~\ref{thm:local-witt}.
\end{proof}

Let $R_{\rm Rees}=\RR[\tau_e:e\in E]$.  The split threshold filtration has one graded line for
each $e$, and its polynomial multi-Rees form is
\begin{equation}\label{eq:polynomial-rees}
 \mathcal Q_A^{\rm Rees}(\boldsymbol\tau;x)
 =\sum_i x_i^2+\sum_e\tau_e^2\ell_e\nu_e(x)^2-(H-1)M(x)^2
 \in R_{\rm Rees}.
\end{equation}
Thus this is an actual polynomial Rees family, not only a collection of numerical weights.  At a
real point $\tau_e=s_e$, put $t_e=s_e^2\geq0$ and write
\begin{equation}\label{eq:local-rees}
 Q_{A,\boldsymbol t}(x)=\sum_i x_i^2+
       \sum_et_e\ell_e\nu_e(x)^2-(H-1)M(x)^2.
\end{equation}
The associated-graded coordinate of $e$ is $z_e=\sqrt{\ell_e}\nu_e$.  Unit specialization
$\tau_e=1$ is $Q_A$; setting $\tau_e=0$ deletes exactly channel $e$.  Conversely every
nonnegative weight vector $\boldsymbol t$ is a real specialization by $s_e=\sqrt{t_e}$.

Choose independent $t_{a,e}$ at all marked points and fuse as in \eqref{eq:fused-form}.

\begin{theorem}[Global multi-Rees Witt law]\label{thm:global-rees}
If all $t_{a,e}\geq0$, the fused Rees form is a positive-definite terminal-plus-graded form plus
the same common-rank square as in Theorem~\ref{thm:global-witt}.  Thus
\[
 \indm(Q_{X,D,\boldsymbol t})+\dim\Rad(Q_{X,D,\boldsymbol t})\leq1,
 \qquad Q_{X,D,\boldsymbol t}>0\quad\text{on }\ker M\setminus\{0\}.
\]
The rank coefficient is independent of all Rees parameters.
\end{theorem}

\begin{proof}
Replace each local positive form by
$\sum_i x_i^2+\sum_et_{a,e}\ell_{a,e}\nu_{a,e}^2$.  Nonnegative parameters preserve all graded
squares and terminal sums remain positive definite.  Fusion changes only common rank.  Repeat the
hyperplane and radical argument in Theorem~\ref{thm:global-witt}.
\end{proof}

\section{Capacity completion and intrinsic phase}\label{sec:capacity}

Write $P_A(x)=x^TPx$ and $M(x)=r^Tx$, $r_i=|i|$.  Then
$Q_A(x)=x^TPx-\lambda(r^Tx)^2$, where $\lambda=H-1$ and $P>0$.  Define
\begin{equation}\label{eq:capacity-data}
                         c=r^TP^{-1}r>0,\qquad \Omega=c^{-1}-\lambda.
\end{equation}

\begin{theorem}[Capacity normal form]\label{thm:capacity}
For $m=r^Tx$, set $z=x-(m/c)P^{-1}r$.  Then $r^Tz=0$, this splitting is unique, and
\begin{equation}\label{eq:capacity-normal}
                         \boxed{Q_A(x)=z^TPz+\Omega m^2.}
\end{equation}
If $\Omega>0$, $Q_A$ is positive definite.  If $\Omega=0$, it is positive semidefinite with
kernel $\RR P^{-1}r$.  If $\Omega<0$, it is nonsingular with exactly one negative eigenvalue.
Moreover
\begin{equation}\label{eq:determinant}
                         \det Q_A=\det(P)c\Omega.
\end{equation}
\end{theorem}

\begin{proof}
The identity $r^TP^{-1}r=c$ gives $r^Tz=0$, uniqueness, and
$z^TP(P^{-1}r)=0$.  Expanding yields \eqref{eq:capacity-normal}; its orthogonal
positive-space-plus-one-line decomposition gives the trichotomy.  The determinant lemma gives
$\det(P-\lambda rr^T)=\det(P)(1-\lambda c)=\det(P)c\Omega$.
\end{proof}

\begin{proposition}[Root-height invariance]\label{prop:root-invariance}
Increasing $H$ by $t>0$ adds the full block $A$ with persistence $t$.  Hence
\[
 P_t=P+trr^T,\qquad \lambda_t=\lambda+t,
 \qquad c_t=\frac{c}{1+tc}.
\]
Thus $c_t^{-1}-\lambda_t=\Omega$ and $\det(P_t)c_t=\det(P)c$.
\end{proposition}

\begin{proof}
The full-block charge is $M=r^Tx$.  Sherman--Morrison gives the formula for $c_t$; the two
invariants follow, using the determinant lemma for the second.
\end{proof}

At puncture $a$, write $c_a,\Omega_a$ for local data and set $\delta_a=2-\Omega_a$.

\begin{theorem}[Euler--capacity refinement]\label{thm:euler-capacity}
Every $x\in V_{X,D}$ has unique centered components $z_a$, and
\begin{equation}\label{eq:euler-capacity}
 \boxed{Q_{X,D}(x)=\sum_a z_a^TP_az_a+\Omega_{X,D}M^2,
 \quad \Omega_{X,D}=\chi(X)+\sum_a(\Omega_a-2)
                    =\chi(X)-\sum_a\delta_a.}
\end{equation}
The coefficient $\Omega_{X,D}$ is the constrained minimum at $M=1$; its sign gives the global
positive, corank-one affine, and Lorentz phases.
\end{theorem}

\begin{proof}
Apply Theorem~\ref{thm:capacity} at every puncture and insert into \eqref{eq:fused-form}.  The rank
coefficient is
$\sum_a\Omega_a+2(1-g-m)=\chi(X)+\sum_a(\Omega_a-2)$.
The centered sum is positive definite and has unique minimum zero.
\end{proof}

\section{Corona elimination as relative vanishing}\label{sec:corona}

Let $X,Y$ be finite-dimensional real vector spaces and
\[
 \mathcal Q(x,y)=
 \begin{pmatrix}x\\y\end{pmatrix}^{T}
 \begin{pmatrix}A&B\\B^T&D\end{pmatrix}
 \begin{pmatrix}x\\y\end{pmatrix},\qquad D>0.
\]
Define
\begin{equation}\label{eq:corona-maps}
 h(x)=(x,-D^{-1}B^Tx),\qquad v(x,y)=y+D^{-1}B^Tx.
\end{equation}

\begin{theorem}[Exact corona--vanishing law]\label{thm:corona}
There is a short exact sequence
\begin{equation}\label{eq:corona-exact}
              0\longrightarrow X\xrightarrow{h}X\oplus Y
              \xrightarrow{v}Y\longrightarrow0.
\end{equation}
Moreover
\begin{equation}\label{eq:schur}
 \boxed{\mathcal Q(x,y)=v(x,y)^TDv(x,y)+x^T(A-BD^{-1}B^T)x.}
\end{equation}
Thus the Schur complement is the unique exact effective boundary energy.  Equality occurs exactly
on $\im h$, and the energy lost under elimination is the positive $D$-norm of the quotient
coordinate $(X\oplus Y)/\im h\cong Y$.
\end{theorem}

\begin{proof}
The map $h$ is injective, $v$ is onto since $v(0,y)=y$, and $v\circ h=0$.  Conversely,
$v(x,y)=0$ implies $(x,y)=h(x)$, proving exactness.  Expansion proves \eqref{eq:schur}; positivity
of $D$ proves the minimization and equality statements.
\end{proof}

\begin{corollary}[Order independence]\label{cor:order}
Successive eliminations of positive corona blocks are independent of elimination order.
\end{corollary}

\begin{proof}
Each step is exact minimization.  Iterated minima over a finite product commute.  Equivalently,
two attained lower bounds for the same full energy bound one another and are equal.
\end{proof}

The residual is not called a vanishing channel merely because its norm is positive: it is the
explicit quotient in \eqref{eq:corona-exact}.  Under a boundary Stokes realization,
Theorem~\ref{thm:transport} carries this quotient to the corresponding relative vanishing object.

\section{Central-port law and arithmetic realization}\label{sec:arithmetic}

This section determines the full scalar kernel of relative Stokes transport and then places an
arithmetic character in that kernel.  The result is stronger than assigning a number to the
negative square in \eqref{eq:intro-law}: the scalar is defined before a divisor is chosen, the
rank functional is forced independently by determinant functoriality, and the $\Xi$ coefficient is
reached by an explicit M\"obius--theta source.

\subsection{The component law for relative channels}

Let $\Gamma$ be the directed interaction graph of a finite Stokes datum: its vertices are the
Stokes circles $i\in I$, and $i\to j$ is an edge whenever the mixed channel
$\Irr\Hom(i,j)$ is nonzero.  Write $V_i=\CC^{r_i}$ and let
$T_i\in\operatorname{GL}(V_i)$ be formal monodromy.
For $\lambda=(\lambda_i)_{i\in I}\in(\CC^\times)^I$, put
$T_i^\lambda=\lambda_iT_i$.  On a mixed Hom channel the twisted transport is
\begin{equation}\label{eq:scalar-hom-transport}
 f\longmapsto T_j^\lambda f(T_i^\lambda)^{-1}
   =\lambda_j\lambda_i^{-1}T_jfT_i^{-1}.
\end{equation}
Let $d_\Gamma\lambda$ denote the edge tuple
$(\lambda_j\lambda_i^{-1})_{i\to j}$.  Components below are components of the underlying
unoriented graph.

\begin{theorem}[Central-component law]\label{thm:central-components}
A scalar tuple is invisible on every nonzero relative Stokes channel if and only if it is constant
on every component of $\Gamma$.  Equivalently,
\begin{equation}\label{eq:central-kernel}
 \ker d_\Gamma=\{\lambda:\lambda_i=\lambda_j
       \text{ whenever }i,j\text{ lie in the same component}\}.
\end{equation}
If $\alpha\in\pi_0(\Gamma)$ and
\[
 M_\alpha(x)=\sum_{i\in\alpha}r_ix_i,
\]
then the scalar $\lambda_\alpha$ on that component acts on the determinant line of
$V(x)=\bigoplus_iV_i^{\oplus x_i}$ with weight $M_\alpha(x)$:
\begin{equation}\label{eq:component-determinant}
 \frac{\det T_x^\lambda}{\det T_x}
       =\prod_{\alpha\in\pi_0(\Gamma)}
          \lambda_\alpha^{M_\alpha(x)}.
\end{equation}
Thus there is one channel-invisible scalar, and one determinant port, per component.
\end{theorem}

\begin{proof}
If $i\to j$ is an edge, invisibility in \eqref{eq:scalar-hom-transport} gives
$\lambda_j\lambda_i^{-1}=1$, hence $\lambda_i=\lambda_j$.  Propagation along an unoriented path,
using equality symmetrically, makes $\lambda$ constant on each component.  Conversely a
componentwise constant tuple cancels in \eqref{eq:scalar-hom-transport}.  This proves
\eqref{eq:central-kernel}.  Finally
\[
 \frac{\det(\lambda_iT_i)^{x_i}}{\det(T_i)^{x_i}}
       =\lambda_i^{r_ix_i}.
\]
Multiplication over $i$ and regrouping by components gives
\eqref{eq:component-determinant}.
\end{proof}

For the coefficient-generated datum of Section~\ref{sec:coefficients}, $\Gamma$ is complete.
Indeed, if $i\ne j$, choose branches $a\in i$ and $b\in j$.  They are distinct, so
$k(a,b)\geq1$ and $h(a,b)>0$; this positive term occurs in the sum defining
$\Irr\Hom(i,j)$.  Hence $\pi_0(\Gamma)$ is a singleton.

\begin{theorem}[Central-port duality]\label{thm:central-port}
For a nonempty coefficient-generated Stokes spectrum, the exact sequence
\begin{equation}\label{eq:central-exact}
 1\longrightarrow\CC^\times\xrightarrow{\Delta}(\CC^\times)^I
   \xrightarrow{d_\Gamma}(\CC^\times)^{E(\Gamma)}
\end{equation}
is exact at $(\CC^\times)^I$.  For every integral dimension vector $x$, restriction of its
determinant character to $\im\Delta$ is
\begin{equation}\label{eq:rank-determinant-character}
              \lambda\longmapsto\lambda^{M(x)},\qquad
              M(x)=\sum_i r_ix_i.
\end{equation}
The same $M$ is the unique port through which a nonpositive term can enter the coarea
polarization.  Thus the metric one-port theorem and the relative-channel kernel theorem have the
same rank functional, obtained by independent constructions.
\end{theorem}

\begin{proof}
Completeness of $\Gamma$ and Theorem~\ref{thm:central-components} identify the kernel with the
diagonal.  Equation \eqref{eq:rank-determinant-character} is
\eqref{eq:component-determinant} for the single component.  Formula
\eqref{eq:intro-law}, or its intrinsic form in Theorem~\ref{thm:coarea}, independently
identifies the only subtractive coordinate as the same ramification-weighted sum $M$.
\end{proof}

\begin{remark}\label{rem:disconnected-ports}
Connectedness is sharp.  If the relative interaction graph has $b$ components, then $b$ independent
central scalars survive, with determinant weights $M_1,\ldots,M_b$.  Consequently no theorem
claiming one universal port can hold for a disconnected Stokes skeleton without an additional
fusion relation.  Coefficient-generated spectra evade this obstruction because their mixed
irregularity graph is complete.
\end{remark}

\subsection{Why arithmetic is neither rank nor formal contact}

Two tempting identifications are impossible.

\begin{proposition}[Rank and contact-stratum no-go]\label{prop:arithmetic-no-go}
Fix a finite-length Stokes category over a connected parameter space.
\begin{enumerate}[label=\textup{(\roman*)}]
\item The continuously varying number
\[
 \delta_n(s)-1=\exp((2\Re s-1)\log n)-1,\qquad n\geq2,
\]
cannot be a physical-rank or Grothendieck-group class.
\item Suppose two equivariant coefficient systems have the same orbit map and, for every
$a,b,k$, the same truth value of $c_{a,k}=c_{b,k}$.  Then they have identical contacts,
threshold blocks, mixed Hom irregularities, Cartan/coarea forms, capacities, and filtered Rees
specializations.  In particular, bare coefficient contact cannot recover the magnitude $\log n$.
\end{enumerate}
\end{proposition}

\begin{proof}
Physical ranks and coordinates of $K_0$ are integral.  A continuous map from a connected space to
this discrete lattice is constant, whereas the displayed function is nonconstant in $\Re s$.

For (ii), $k(a,b)$ is the largest exponent at which coefficient equality fails.  Equality of all
these truth values therefore preserves every $k(a,b)$ and $h(a,b)$, hence every threshold
partition and persistence length.  The definitions in Sections~\ref{sec:coefficients} and
\ref{sec:rees} then preserve all listed quantities.
\end{proof}

The proposition forces any continuously varying, relative-channel-invisible arithmetic datum into
monodromy or extension data over a fixed irregular skeleton.  Theorem~\ref{thm:central-port} then
forces a scalar such datum into the common center.  This is a classification, not a choice of
convenient coordinates.

\subsection{The pre-divisor arithmetic character}

For $n\geq1$ and $s\in\CC$, define
\begin{equation}\label{eq:delta-character}
 \delta_n(s)=\exp((2\Re s-1)\log n)>0,\qquad
 \chi_{n,s}(k)=\delta_n(s)^k\quad(k\in\pi_1(S^1)=\mathbb Z).
\end{equation}
Let $\mathcal L_{n,s}$ be the rank-one Betti local system represented by $\chi_{n,s}$.  Tensoring
every Stokes block with $\mathcal L_{n,s}$ replaces $T_i$ by $\delta_n(s)T_i$.  This family is
real analytic, not holomorphic, in $s$; no holomorphic parameter-space assertion is used.

\begin{theorem}[Arithmetic realization of the central port]\label{thm:arithmetic-center}
The family $\mathcal L_{n,s}$ is defined before any zero or divisor is selected and satisfies:
\begin{enumerate}[label=\textup{(\roman*)}]
\item it fixes all exponential differences, analytic Stokes walls, mixed Hom transports, directed
microstalk ranks, coefficient cuts, coarea forms, capacities, corona quotients, and Rees
specializations;
\item on the determinant line of $V(x)$ its relative holonomy is
\begin{equation}\label{eq:det-holonomy}
                       \delta_n(s)^{M(x)};
\end{equation}
\item for $m,n\geq1$,
\begin{equation}\label{eq:delta-multiplicative}
 \delta_{mn}(s)=\delta_m(s)\delta_n(s),\qquad
 \delta_{n^{M(x)}}(s)=\delta_n(s)^{M(x)};
\end{equation}
\item critical reflection dualizes the local system:
\begin{equation}\label{eq:delta-reflection}
                  \delta_n(1-\overline s)=\delta_n(s)^{-1};
\end{equation}
\item among scalar twists invisible on every mixed channel, it is uniquely determined by its
holonomy on any one Stokes block.
\end{enumerate}
For $n\geq2$, $\mathcal L_{n,s}$ is unitary if and only if $\Re s=1/2$.
\end{theorem}

\begin{proof}
The tensor factor is regular singular and common to every block.  It changes no irregular
exponential factor, so all difference, wall, contact, microlocal, and filtered constructions are
unchanged; cancellation on mixed channels is also immediate from
\eqref{eq:scalar-hom-transport}.  Equation \eqref{eq:det-holonomy} is
Theorem~\ref{thm:central-port}.  The first identity in \eqref{eq:delta-multiplicative} follows from
$\log(mn)=\log m+\log n$; iteration gives the second.  Since
$2\Re(1-\overline s)-1=-(2\Re s-1)$, \eqref{eq:delta-reflection} follows.  Uniqueness is
Theorem~\ref{thm:central-components}.  Finally a positive real rank-one holonomy is unitary exactly
when it is one; for $n\geq2$, positivity of $\log n$ makes this equivalent to
$\Re s=1/2$.
\end{proof}

\subsection{A target-free source for all completed-\texorpdfstring{$\Xi$}{Xi} jets}

We now show that the $\Xi$ coefficient used below is reached by arithmetic source data rather than
inserted at a chosen zero.  Set
\[
 \xi(s)=\frac12s(s-1)\pi^{-s/2}\Gamma(s/2)\zeta(s),\qquad
 \mathcal X(w)=\xi\left(\frac12+iw\right).
\]
Define
\[
 h(t)=2\pi(2\pi t^4-3t^2)e^{-\pi t^2},\qquad
 \Theta(x)=\sum_{\ell\geq0}e^{x/2}h(\ell e^x).
\]
This is the even theta-cusp Schwartz kernel, normalized so that
\begin{equation}\label{eq:theta-xi-transform}
             \int_{\RR}e^{-iwx}\Theta(x)\,dx=\mathcal X(w)
             \qquad(|\Im w|<1).
\end{equation}
The identity follows by termwise Mellin integration of the polynomial Gaussian and the theta
functional equation; this normalization is the standard Fourier form of the completed zeta
function \cite{Titchmarsh1986}.

For $m\geq1$, let
\[
 a_m(N)=
 \begin{cases}
   \mu(q),&N=q^m,\\
   0,&\text{otherwise},
 \end{cases}
\]
and, for each jet grade $k\geq0$, form the source
\begin{equation}\label{eq:mobius-theta-source}
 \Psi_{m,k}(x)=\sum_{N\geq1}
    a_m(N)N^{-1/2}\bigl(-i(x-\log N)\bigr)^k\Theta(x-\log N).
\end{equation}
Every object in \eqref{eq:mobius-theta-source} is defined without a spectral point.

\begin{proposition}[Source-to-$\Xi$ jet landing]\label{prop:source-xi-jet}
Put $s_w=\frac12+iw$.  If $|\Im w|<1$ and $\Re s_w>1/m$, then
\begin{equation}\label{eq:source-xi-transform}
 \int_{\RR}e^{-iwx}\Psi_{m,k}(x)\,dx
      =\zeta(ms_w)^{-1}\mathcal X^{(k)}(w).
\end{equation}
For fixed $w$ with $\Re s_w>0$, the left side tends cofinally, as $m\to\infty$, to
$\mathcal X^{(k)}(w)$.
\end{proposition}

\begin{proof}
The strip bound makes every polynomially weighted theta packet integrable.  The half-plane
condition gives
$\sum_N|a_m(N)|N^{-\Re s_w}<\infty$, so Fubini applies.  In the $N$th term set
$y=x-\log N$.  Differentiating \eqref{eq:theta-xi-transform} under the integral gives
\[
 \int_{\RR}e^{-iwy}(-iy)^k\Theta(y)\,dy=\mathcal X^{(k)}(w).
\]
Translation contributes $N^{-iw}$, and therefore the coefficient sum is
\[
 \sum_{N\geq1}a_m(N)N^{-s_w}
   =\sum_{q\geq1}\mu(q)q^{-ms_w}
   =\zeta(ms_w)^{-1}.
\]
This proves \eqref{eq:source-xi-transform}.  If $\Re s_w>0$, then
$\zeta(ms_w)\to1$: the absolutely convergent tail
$\sum_{q\geq2}q^{-m\Re s_w}$ tends to zero.  The stated cofinal limit follows.
\end{proof}

Let $\rho$ be a nontrivial divisor point of $\xi$ and put
\begin{equation}\label{eq:xi-leading-jet}
 w_\rho=-i(\rho-\tfrac12),\qquad
 r_\rho=\operatorname{ord}_{w_\rho}\mathcal X,\qquad
 c_\rho=\mathcal X^{(r_\rho)}(w_\rho)\ne0.
\end{equation}
Since $0<\Re\rho<1$, $|\Im w_\rho|<1$ and $\Re s_{w_\rho}>0$.  Thus
Proposition~\ref{prop:source-xi-jet} reaches every jet at $w_\rho$, all lower grades vanish, and
the first nonzero grade lands at $c_\rho$ without division by $\mathcal X$.

\subsection{Rational-adelic trace and the rank-powered transgression}

For completeness we construct the cyclic source used for the final landing.  Let
\[
 \mathcal B=\mathcal S(\RR)\otimes
     \mathcal B_c(\mathbb A_{\QQ,f}),\qquad G=\QQ^{\times},
\]
where the second factor is the algebraic Bruhat space on finite adeles.  The algebraic crossed
product $A_0=\mathcal B\rtimes_{\rm alg}G$ consists of finite sums $aU_q$ with
\begin{equation}\label{eq:crossed-product}
                  (aU_q)(bU_\ell)=a(q\cdot b)U_{q\ell}.
\end{equation}
Let $A_{\rm ad}$ be its minimal unitization.  Evaluation at the zero adele is multiplicative on
$\mathcal B$ and invariant under scaling.  Hence
\begin{equation}\label{eq:epsilon-zero}
 \varepsilon_0(\alpha,F)=\alpha+\operatorname{ev}_0(F(1))
\end{equation}
is a trace on $A_{\rm ad}$.  Indeed, on two crossed-product monomials both traces vanish unless
$q\ell=1$; in that case scaling invariance and commutativity in $\mathcal B$ give
$\operatorname{ev}_0(a(q\cdot b))=\operatorname{ev}_0(b(\ell\cdot a))$.
The unitization terms are immediate.

Take the genuine degree-one cyclic bar
\[
 C_1(A_{\rm ad})=A_{\rm ad}\otimes A_{\rm ad},\qquad
 \tau(a_0\otimes a_1)=a_1\otimes a_0,
\]
and its cyclic trace
\begin{equation}\label{eq:cyclic-trace}
 \varepsilon_0^{(1)}(a_0\otimes a_1)
     =\varepsilon_0(a_0a_1),\qquad
 \varepsilon_0^{(1)}\tau=\varepsilon_0^{(1)}.
\end{equation}
Let $g=e^{-\pi x^2}\otimes\mathbf 1_{\widehat{\mathbb Z}}$, so
$\operatorname{ev}_0(g)=1$, and for an ordinary prime $p$ define
\begin{equation}\label{eq:prime-cycle}
                u_p=(gU_p)\otimes(gU_{p^{-1}}).
\end{equation}

\begin{lemma}[Trace-one prime cycle]\label{lem:prime-trace-one}
For every ordinary prime $p$,
\[
                         \varepsilon_0^{(1)}(u_p)=1.
\]
\end{lemma}

\begin{proof}
The scaling labels in \eqref{eq:prime-cycle} multiply to one.  Since zero is fixed by scaling,
\[
 \varepsilon_0^{(1)}(u_p)
 =\operatorname{ev}_0\bigl(g(p\cdot g)\bigr)
 =\operatorname{ev}_0(g)^2=1.
\]
\end{proof}

Define
\[
 C(y)=(\varepsilon_0^{(1)}(y),0),\qquad
 \operatorname{Res}(z_0,z_1)=z_0+z_1,
\]
and, for every independent grade $k$ and spectral parameter $w$,
\[
 L_{k,w}(y)=\mathcal X^{(k)}(w)C(y).
\]
Proposition~\ref{prop:source-xi-jet} realizes $L_{k,w}$ as the cofinal limit of the same expression
with $\mathcal X^{(k)}(w)$ replaced by the transform of $\Psi_{m,k}$.  For a positive integer
clock $q$, put
\begin{equation}\label{eq:paracyclic-rotation}
 R_{q,s}(z_0,z_1)=(z_1,\delta_q(s)z_0),\qquad
 D_{q,s,k,w}=R_{q,s}L_{k,w}-L_{k,w}\tau.
\end{equation}
Notice that $R_{q,s}^2=\delta_q(s)\operatorname{id}$; this is a paracyclic, not generally cyclic,
target rotation.

\begin{theorem}[Rank-port/$\Xi$ determinant law]\label{thm:rank-port-xi}
Let $x\in\NN^I$, $n\geq1$, and let $\rho$ be any nontrivial divisor point.  Then, for every
$y\in C_1(A_{\rm ad})$,
\begin{equation}\label{eq:rank-port-xi}
 \boxed{\;
 \operatorname{Res}\bigl(D_{n^{M(x)},\rho,r_\rho,w_\rho}y\bigr)
 =\bigl(\delta_n(\rho)^{M(x)}-1\bigr)c_\rho
      \varepsilon_0^{(1)}(y).
 \;}
\end{equation}
The multiplier on the right is precisely holonomy minus identity on the determinant line of
$V(x)\otimes\mathcal L_{n,\rho}$.  On the trace-one prime cycle,
\begin{equation}\label{eq:prime-rank-port-xi}
 \operatorname{Res}\bigl(D_{n^{M(x)},\rho,r_\rho,w_\rho}u_p\bigr)
       =\bigl(\delta_n(\rho)^{M(x)}-1\bigr)c_\rho.
\end{equation}
If $n\geq2$ and $M(x)>0$, this quantity vanishes if and only if
$\Re\rho=1/2$.
\end{theorem}

\begin{proof}
Write $a=c_\rho\varepsilon_0^{(1)}(y)$.  Cyclic trace invariance gives
\[
 L_{r_\rho,w_\rho}(\tau y)=(a,0),\qquad
 R_{q,\rho}L_{r_\rho,w_\rho}(y)=(0,\delta_q(\rho)a).
\]
Subtraction and coordinate sum yield
\[
 \operatorname{Res}(D_{q,\rho,r_\rho,w_\rho}y)
       =(\delta_q(\rho)-1)a.
\]
Take $q=n^{M(x)}$ and use \eqref{eq:delta-multiplicative}.  By
\eqref{eq:det-holonomy}, $\delta_n(\rho)^{M(x)}$ is exactly the determinant-line holonomy ratio.
Lemma~\ref{lem:prime-trace-one} gives \eqref{eq:prime-rank-port-xi}.  Since $c_\rho\ne0$,
$n\geq2$, and $M(x)>0$, its vanishing is equivalent to
\[
 \exp\bigl(M(x)(2\Re\rho-1)\log n\bigr)=1.
\]
The exponent is real and $M(x)\log n>0$, so this holds exactly when $\Re\rho=1/2$.
\end{proof}

The theorem gives the following arithmetic defect diagram:
\[
\begin{array}{ccc}
 C_1(A_{\rm ad})&\xrightarrow{\quad L_{r_\rho,w_\rho}\quad}&\CC^2\\
 \big\downarrow{\tau}&&\big\downarrow{R_{n^{M(x)},\rho}}\\
 C_1(A_{\rm ad})&\xrightarrow{\quad L_{r_\rho,w_\rho}\quad}&\CC^2
 \xrightarrow{\ \operatorname{Res}\ }\CC .
\end{array}
\]
Its failure to commute is measured by determinant holonomy minus identity, and the exponent is
the independently forced Stokes rank port $M(x)$.

\subsection{Reflection-paired vanishing and the Stokes--Mertens law}

Put $s^\vee=1-\overline s$.  At a nontrivial divisor $\rho$, the functional equation and reality
of $\xi$ make $\rho^\vee$ a divisor of the same multiplicity.  Fix a positive integer $M$ and,
for a prime $p$, normalize the two rank-powered transgressions by their nonzero leading jets:
\begin{equation}\label{eq:normalized-reflected-defects}
 \begin{aligned}
 a_{p,M}(\rho)
  &=\frac{\operatorname{Res}(D_{p^M,\rho,r_\rho,w_\rho}u_p)}{c_\rho}
    =\delta_p(\rho)^M-1,\\
 a_{p,M}(\rho^\vee)
  &=\frac{\operatorname{Res}
      (D_{p^M,\rho^\vee,r_{\rho^\vee},w_{\rho^\vee}}u_p)}{c_{\rho^\vee}}
    =\delta_p(\rho)^{-M}-1.
 \end{aligned}
\end{equation}
The division here removes only the source jet, which is nonzero by definition; it does not alter
or postulate a divisor.  For arbitrary $s\in\CC$ define the same pre-divisor defects by
$a_{p,M}(s)=\delta_p(s)^M-1$ and set
\begin{equation}\label{eq:stokes-prime-energy}
 \mathcal E_{N,M}(s)
   =-\sum_{p\leq N}\frac{a_{p,M}(s)a_{p,M}(s^\vee)}p.
\end{equation}

\begin{theorem}[Stokes--Mertens reflection and vanishing law]
\label{thm:stokes-mertens}
Let $M>0$.
\begin{enumerate}[label=\textup{(\roman*)}]
\item Reflection pairing is the exact positive factorization
\begin{equation}\label{eq:stokes-energy-hyperbolic}
 \boxed{\quad
 \mathcal E_{N,M}(s)
   =\sum_{p\leq N}
      \frac{4\sinh^2(M(\Re s-\tfrac12)\log p)}p\geq0.
 \quad}
\end{equation}
At a divisor this is the negative product of the two normalized rational-adelic $\Xi$
transgressions in \eqref{eq:normalized-reflected-defects}.
\item For every $N\geq2$,
\begin{equation}\label{eq:finite-stokes-vanishing}
             \mathcal E_{N,M}(s)=0
             \quad\Longleftrightarrow\quad \Re s=\frac12.
\end{equation}
Thus one prime already gives an exact arithmetic vanishing locus, while the full prime system
quantifies transverse escape.
\item Let $s_N\in\CC$ be arbitrary and put
\(
 c_N=M(\Re s_N-\tfrac12)\log N.
\)
Then
\begin{equation}\label{eq:stokes-mertens-trichotomy}
 \begin{array}{rcl}
 c_N\longrightarrow0
   &\Longrightarrow&\mathcal E_{N,M}(s_N)\longrightarrow0,\\[2mm]
 c_N\longrightarrow c\in\RR
   &\Longrightarrow&\displaystyle
      \mathcal E_{N,M}(s_N)\longrightarrow
      \Phi(c):=4\int_0^1\frac{\sinh^2(cu)}u\,du,\\[3mm]
 |c_N|\longrightarrow\infty
   &\Longrightarrow&\mathcal E_{N,M}(s_N)\longrightarrow\infty.
 \end{array}
\end{equation}
The middle profile is even, and $\Phi(c)>0$ for $c\ne0$.  In particular the critical-window
limit is independent of every Stokes multiplicity and contact coefficient: their sole surviving
quantity $M$ has been absorbed by the forced scaling.
\end{enumerate}
\end{theorem}

\begin{proof}
Write $g=\delta_p(s)^M>0$.  Reflection gives
$\delta_p(s^\vee)^M=g^{-1}$, and hence
\[
 -a_{p,M}(s)a_{p,M}(s^\vee)
 =-(g-1)(g^{-1}-1)=g+g^{-1}-2.
\]
Since
$g=\exp(2M(\Re s-\tfrac12)\log p)$, the last expression is
$4\sinh^2(M(\Re s-\tfrac12)\log p)$.  Summing proves
\eqref{eq:stokes-energy-hyperbolic}; at a divisor,
\eqref{eq:normalized-reflected-defects} follows from
Theorem~\ref{thm:rank-port-xi}.  If $N\geq2$, the sum contains $p=2$ and consists of nonnegative
terms.  It is zero exactly when
$M(\Re s-\tfrac12)\log2=0$, which proves \eqref{eq:finite-stokes-vanishing}.

For completeness we prove the limiting assertion from the prime form of Mertens's first theorem
\cite{Mertens1874}.
Put
\[
 S_N=\sum_{p\leq N}\frac{(\log p)^2}{p}.
\]
Partial summation applied to
$\sum_{p\leq x}(\log p)/p=\log x+O(1)$ gives
\begin{equation}\label{eq:prime-susceptibility}
                 S_N=\frac12(\log N)^2+O(\log N).
\end{equation}
Indeed, if $A(x)=\sum_{p\leq x}(\log p)/p$, then
$S_N=(\log N)A(N)-\int_2^NA(t)\,dt/t+O(1)$, and substitution of
$A(t)=\log t+O(1)$ yields \eqref{eq:prime-susceptibility}.

Write
\[
 E_N^{\rm pr}(\varepsilon)=
   \sum_{p\leq N}\frac{4\sinh^2(\varepsilon\log p)}p,
 \qquad
 \mathcal E_{N,M}(s)=E_N^{\rm pr}\bigl(M(\Re s-\tfrac12)\bigr).
\]
Set $\varepsilon_N=M(\Re s_N-\tfrac12)$, so $c_N=\varepsilon_N\log N$.
The inequalities $|u|\leq|\sinh u|$ and
$|\sinh u|\leq|u|\cosh C$ for $|u|\leq C$ give
\begin{equation}\label{eq:stokes-energy-sandwich}
 4\varepsilon_N^2S_N\leq\mathcal E_{N,M}(s_N),
 \qquad
 \mathcal E_{N,M}(s_N)\leq4\cosh^2(C)\varepsilon_N^2S_N
 \quad(|c_N|\leq C).
\end{equation}
Together with \eqref{eq:prime-susceptibility}, these bounds prove the first and third regimes.

For the middle regime first take $c_N=c$.  The finite measures
\[
 \nu_N=\frac1{\log N}\sum_{p\leq N}\frac{\log p}{p}
             \,\delta_{\log p/\log N}
\]
converge weakly to Lebesgue measure on $[0,1]$, again by Mertens's theorem and partial summation.
The function $u\mapsto4\sinh^2(cu)/u$, continuously extended by zero at $u=0$, is continuous;
therefore
\[
 \mathcal E_{N,M}(s_N)=
 \int_0^1\frac{4\sinh^2(cu)}u\,d\nu_N(u)\longrightarrow\Phi(c).
\]
If merely $c_N\to c$, all $c_N$ lie in one compact interval eventually.  The mean-value theorem
and \eqref{eq:prime-susceptibility} give, for a constant $K$ depending only on that interval,
\[
 \left|E_N^{\rm pr}(c_N/\log N)-E_N^{\rm pr}(c/\log N)\right|
 \leq K|c_N-c|\frac{S_N}{(\log N)^2}\longrightarrow0.
\]
This proves the middle regime.  Evenness and strict positivity of $\Phi$ away from zero are
immediate from its integrand.
\end{proof}

The theorem turns the central representation law into a number-theoretic scaling theorem and a
literal finite vanishing theorem.  It does not force the energy to vanish on every $\Xi$ divisor.
For any fixed $n,N\geq2$ and any positive-rank $x$, it sharpens the classifier to
\begin{equation}\label{eq:rh-unitarity-form}
\begin{aligned}
 \mathrm{RH}
 &\Longleftrightarrow
   \mathcal L_{n,\rho}\ \text{is unitary for every nontrivial divisor point }\rho\,,\\
 &\Longleftrightarrow
   \operatorname{Res}(D_{n^{M(x)},\rho,r_\rho,w_\rho}u_p)=0
   \quad\text{for every such }\rho\text{ and every prime }p,\\
 &\Longleftrightarrow
   \mathcal E_{N,M(x)}(\rho)=0
   \quad\text{for every such }\rho.
\end{aligned}
\end{equation}
Equation \eqref{eq:rh-unitarity-form} is an exact classifier, not a proof of its left side.
The construction supplies the carrier and identifies its only possible invisible port.  What
remains is an independent arithmetic theorem forcing global triviality, equivalently
source-null-homology of the displayed paracyclic defect.  No positivity, reflection symmetry,
or Stokes realization established here supplies that missing annihilation.

\section{The Stokes deformation complex}\label{sec:deformation}

We derive the dimension identity from tangent data.  Fix $p$ punctures and integers
\[
 g,\quad G=\dim\mathfrak g,\quad z=\dim Z(\mathfrak g),\quad
 h_a,\ s_a,\ k_a\quad(1\leq a\leq p),
\]
where $h_a$ is formal-centralizer dimension, $s_a$ directed Stokes-root dimension, and $k_a$
prescribed conjugacy-class dimension.  Set
\begin{align}\label{eq:factor-counts}
 A_1&=2gG+\sum_a(G+h_a+s_a+k_a),\\
 R&=G-z+\sum_ah_a.\notag
\end{align}
At a smooth stable reduction point, let
\[
 d_{\rm act}:\CC^R\to\CC^{A_1},\qquad d_\mu:\CC^{A_1}\to\CC^R
\]
be effective gauge and moment derivatives.  Smooth stability here means $d_{\rm act}$ is
injective, $d_\mu$ is surjective, and $d_\mu d_{\rm act}=0$.

\begin{theorem}[Source-derived tangent dimension]\label{thm:tangent}
The reduced tangent $T=\ker d_\mu/\im d_{\rm act}$ satisfies
\begin{equation}\label{eq:tangent}
 \boxed{\dim T=2z+(2g-2+p)G+\sum_a(s_a+k_a-h_a).}
\end{equation}
\end{theorem}

\begin{proof}
Surjectivity gives $\dim\ker d_\mu=A_1-R$; injectivity gives
$\dim\im d_{\rm act}=R$.  Thus $\dim T=A_1-2R$.  Substitute \eqref{eq:factor-counts}.
\end{proof}

Form the three-term complex
\begin{equation}\label{eq:def-complex}
 \CC^z\oplus\CC^R\xrightarrow{d_0}\CC^{A_1}
 \xrightarrow{d_1}\CC^R\oplus\CC^z,
 \quad d_0(c,u)=d_{\rm act}u,\quad d_1(v)=(d_\mu v,0).
\end{equation}

\begin{theorem}[Self-dual Euler count]\label{thm:euler-count}
Sequence \eqref{eq:def-complex} is a complex and
$\dim H^0=z$, $H^1\cong T$, $\dim H^2=z$.  Hence
\begin{equation}\label{eq:euler-count}
 \boxed{\chi(\Def_{\rm St})=(2-2g-p)G+\sum_a(h_a-s_a-k_a).}
\end{equation}
\end{theorem}

\begin{proof}
The complex identity is $d_\mu d_{\rm act}=0$.  Injectivity leaves exactly the central summand in
$H^0$; degree-one cohomology is $T$; surjectivity leaves exactly the final central summand in
$H^2$.  Alternate dimensions and use Theorem~\ref{thm:tangent}.
\end{proof}

By Theorem~\ref{thm:micro-index}, $s_a$ is an integer directed microstalk sum.  Thus (M), (D), and
(E) give a closed chain from coefficient differences through vanishing quotients to tangent
dimension.  The rank--nullity calculation independently recovers the quasi-Hamiltonian polynomial
in Proposition~\ref{prop:qh-fusion}; agreement is a compatibility check, not a second definition.

\section{Type-A marking legs}\label{sec:legs}

Attach a path of length $l_i\geq0$ to orbit vertex $i$.  Put $y_{i,0}=x_i$,
$y_{i,l_i+1}=0$, and $d_{i,k}=y_{i,k}-y_{i,k+1}$.  Then $\sum_kd_{i,k}=x_i$.

\begin{lemma}[Path completion]\label{lem:path}
If $d_0+\cdots+d_l=x$, then
\begin{equation}\label{eq:path}
 \sum_{k=0}^{l}d_k^2=
 \sum_{k=0}^{l}\left(d_k-\frac{x}{l+1}\right)^2+\frac{x^2}{l+1}.
\end{equation}
\end{lemma}

\begin{proof}
Expand and use $\sum_kd_k=x$.
\end{proof}

\begin{theorem}[Leg-stable Witt law]\label{thm:legs}
Replacing each terminal square $x_i^2$ by its full type-$A$ path energy gives
\begin{align}\label{eq:leg-law}
 \widehat Q_A
 &=\sum_i\sum_{k=0}^{l_i}\left(d_{i,k}-\frac{x_i}{l_i+1}\right)^2
   +\sum_i\frac{x_i^2}{l_i+1}
   +\sum_e\ell_e\nu_e(x)^2-(H-1)M(x)^2.
\end{align}
Thus finite marking legs do not create a second negative or null direction.
\end{theorem}

\begin{proof}
Apply Lemma~\ref{lem:path} to every leg.  Fluctuation squares form a positive complement.  The
remaining positive matrix has diagonal $\operatorname{diag}((l_i+1)^{-1})>0$ plus the cut Gram
matrix and is perturbed by the same one rank square.  Repeat Theorem~\ref{thm:global-witt}.
\end{proof}

\begin{corollary}[Nonabelian Hodge diagram law]\label{cor:nah-diagram}
Let $\widehat C$ be the symmetric generalized Cartan matrix of a genus-zero nonabelian Hodge
diagram obtained from finite ramified Stokes data with its type-$A$ marking legs.  Then
\begin{equation}\label{eq:nah-diagram-law}
                  \indm(\widehat C)+\dim\ker\widehat C\leq1.
\end{equation}
Thus every such diagram is finite, affine of corank one, or Lorentzian.  The assertion is invariant
under every relabelling that preserves its generalized Cartan matrix.
\end{corollary}

\begin{proof}
Apply Proposition~\ref{prop:published-Cartan} to the cyclic-continuation orbits of the defining
finite Puiseux spectrum and Theorem~\ref{thm:legs} to its marking legs.  Relabelling preserves
matrix inertia.
\end{proof}

\section{Spectral consequence and sharpness}\label{sec:spectral}

Suppose a nonabelian Hodge diagram is an honest finite graph with symmetric adjacency matrix $B$,
and let $\lambda_1(B)\geq\cdots\geq\lambda_n(B)$.

\begin{theorem}[Second-eigenvalue obstruction]\label{thm:lambda-two}
If $n\geq2$, then $\lambda_2(B)\leq2$.  If $\lambda_1(B)=2$, eigenvalue $2$ is simple.  If
$\lambda_1(B)>2$, exactly one eigenvalue exceeds $2$ and none equals $2$.
\end{theorem}

\begin{proof}
Eigenvalues of $C$ are $2-\lambda_i(B)$.  Apply \eqref{eq:nah-diagram-law}.
\end{proof}

\begin{example}[Dou\c{c}ot's twisted triangle; sharpness]\label{ex:triangle}
Dou\c{c}ot obtains the graph below from the three Stokes circles
$\langle z^3\rangle$, $\langle z^{4/3}\rangle$, and $\langle z^{3/2}\rangle$
\cite[Proposition~3.3]{Doucot2025}.  The example and its edge multiplicities are his; the point
here is the sharp inertia consequence of the coarea law.  For
\[
 B=\begin{pmatrix}0&6&4\\6&0&3\\4&3&0\end{pmatrix},\qquad C=2I-B,
\]
one has $\one^TC\one=-20$, $\det C=-258$, and $\operatorname{tr}C=6$.  Hence inertia is
$(2,1,0)$.  Explicitly, with
$u=x-3y-2z$ and $v=y+\frac{15}{16}z$,
\[
 (x,y,z)C(x,y,z)^T=2u^2-16v^2+\frac{129}{16}z^2.
\]
The triangular coordinate change has determinant one, so this is a congruence diagonalization.
One negative direction is sharp.
\end{example}

\section{Novelty boundary and falsifiers}\label{sec:novelty}

Finite Stokes filtrations and categorical forms are classical \cite{Sabbah}; irregular
Riemann--Hilbert theory supplies analytic realizations \cite{DKRH}; ultrametric matrices have tree
and inverse $M$-matrix theories \cite{Martinez1994,McDonald1995,Nabben1995}; fission and wild
character spaces are developed in \cite{Boalch2014}; fission trees and ultrametric graph constraints
appear in \cite{Boalch2025,Doucot2025}; and Hiroe gives a scalar spectral-curve Euler identity
\cite{Hiroe2019}.  The diagram formula \eqref{eq:published-B} and the dimension identity
$\dim\Mcal_B=2-(d,d)$ are due to Boalch--Yamakawa and Dou\c{c}ot
\cite{BoalchYamakawa2020,Doucot2021}.  These are inputs, not novelty claims.
Rank-one local systems, determinant scaling, the adeles-class-space source, and the classical
theta transform of $\xi$ are likewise not claimed as new separately
\cite{Connes1999,ConnesConsaniMarcolli2007,Titchmarsh1986}.

New claim is the complete polarization and its consequences:
\begin{align*}
 \text{Puiseux coefficients}&\Longrightarrow\text{canonical cuts}
 \Longrightarrow\text{positive coarea minus one rank line}\\
 &\Longleftrightarrow\text{published Cartan form}
 \Longrightarrow \indm+\dim\Rad\leq1.
\end{align*}
This does not follow from one scalar evaluation at a dimension vector.  It controls the entire
real marked lattice, identifies its positive hyperplane and affine ray, and excludes all
higher-signature Kac--Moody types among realized diagrams.  Globally, all punctures and handles
share the same rank port, yielding the arbitrary-genus Euler--capacity law.  Analytic wall and
microstalk calculations explain why mixed Hom irregularity is the directed local multiplicity
entering that polarization.

Arithmetic novelty is the exact bridge, not its elementary pieces.  The universal component law
classifies every relative-channel-invisible scalar and shows that its determinant weights are
component physical ranks.  Coefficient generation collapses these to the same single functional
$M$ forced by the coarea theorem.  The pre-divisor M\"obius--theta and rational-adelic source then
lands the rank-powered determinant defect at every completed-$\Xi$ divisor point as
$(\delta_n(\rho)^{M(x)}-1)c_\rho$.  No zero position enters construction of source, character,
interaction graph, or determinant exponent.
Reflection pairing before codiagonal erasure is an additional theorem-level consequence: the
negative product of the two normalized unsquared transgressions is exactly the positive energy
\eqref{eq:stokes-energy-hyperbolic}.  Mertens's theorem then gives the three-regime limit
\eqref{eq:stokes-mertens-trichotomy}.  Neither Mertens's estimate nor a hyperbolic-square identity
is claimed separately as new; candidate novelty is that the independently forced Stokes
determinant port is precisely the renormalization exponent and that the rational-adelic $\Xi$
defects furnish the two reflected factors.

Any of the following would falsify the theorem as stated.
\begin{enumerate}[label=(\arabic*),leftmargin=2.2em]
\item an equivariant spectrum violating Lemma~\ref{lem:three-point};
\item failure of finite coarea identity \eqref{eq:coarea};
\item a zero of $\operatorname{Re}(ce^{-ik\theta})$ outside
      \eqref{eq:analytic-wall-labels}, a repeated wall, or failure of the transverse-sign law;
\item a directed exponential germ whose quotient microstalk mass differs from its degree;
\item a nonzero compatible rank-zero charge of nonpositive fused energy;
\item a nonnegative Rees specialization whose negative index plus nullity exceeds one;
\item a positive corona block violating \eqref{eq:schur} or \eqref{eq:corona-exact};
\item a nonconstant scalar twist invisible on a connected interaction graph;
\item determinant response of a common scalar differing from $\lambda^{M(x)}$;
\item failure of the trace-one identity or of the rank-powered transgression
      \eqref{eq:rank-port-xi};
\item failure of reflection factorization \eqref{eq:stokes-energy-hyperbolic}, finite vanishing
      \eqref{eq:finite-stokes-vanishing}, or any regime in
      \eqref{eq:stokes-mertens-trichotomy};
\item a realized nonabelian Hodge graph with $\lambda_2(B)>2$.
\end{enumerate}
Conversely, this is not an analytic existence criterion: a quadratic form satisfying the Witt law
need not arise from a meromorphic connection.

\section{Conclusion}

The universal law identifies the published generalized Cartan polarization with a positive Stokes
coarea space and one physical-rank line.  Principal conormals, directed microlocal Morse ranks,
filtered graded channels, corona residuals, marking legs, and local terms of the deformation
complex are exact finite presentations of its positive part.  Fusion and topology cannot create
another defect because they see local data only through common physical rank.

This single structure explains analytic wall count and orientation, local Lorentz bounds, affine
simplicity, capacity walls, handle
monotonicity, Rees stability, Schur order independence, and the adjacency bound
$\lambda_2\leq2$.  It also replaces the former interface between coefficient contacts and
microlocal multiplicities by an exact categorical quotient calculation.

The central-component law adds an arithmetic meaning to the same port.  Relative Stokes channels
see only scalar ratios, so connected coefficient spectra leave one invisible center; determinant
functoriality reads that center with exponent $M$.  The half-density character and explicit
rational-adelic prime cycle then convert its determinant holonomy defect into the first nonzero
$\Xi$ jet, with exact value $(\delta_n(\rho)^{M(x)}-1)c_\rho$.  This supplies a new proof
architecture for any arithmetic selection theorem.  Reflection now turns two such unsquared
defects into a prime energy with an exact finite vanishing locus and a universal rank-scaled
Mertens profile.  This quantifies the one channel but does not annihilate it: a proof of RH must
still kill the named central Bockstein channel, not reconstruct microsupport or add another
positive form.  Establishing that global source-null-homology remains equivalent to RH and is not
proved here.

\begin{thebibliography}{99}
\bibitem{Boalch2014}
P.~Boalch, \emph{Geometry and braiding of Stokes data; fission and wild character varieties},
Ann. of Math. (2) \textbf{179} (2014), 301--365.
\bibitem{Boalch2025}
P.~Boalch, \emph{Counting the fission trees and nonabelian Hodge graphs (untwisted case)},
J. Geom. Phys. \textbf{214} (2025), 105512.
\bibitem{BoalchYamakawa2015}
P.~Boalch and D.~Yamakawa, \emph{Twisted wild character varieties},
arXiv:1512.08091.
\bibitem{BoalchYamakawa2020}
P.~Boalch and D.~Yamakawa, \emph{Diagrams for nonabelian Hodge spaces on the affine line},
C. R. Math. \textbf{358} (2020), 59--65.
\bibitem{Connes1999}
A.~Connes, \emph{Trace formula in noncommutative geometry and the zeros of the Riemann zeta
function}, Selecta Math. (N.S.) \textbf{5} (1999), 29--106.
\bibitem{ConnesConsaniMarcolli2007}
A.~Connes, C.~Consani and M.~Marcolli, \emph{The Weil proof and the geometry of the adeles
class space}, arXiv:math/0703392.
\bibitem{DKRH}
A.~D'Agnolo and M.~Kashiwara, \emph{Riemann--Hilbert correspondence for holonomic
$\mathcal D$-modules}, Publ. Math. Inst. Hautes Etudes Sci. \textbf{123} (2016), 69--197.
\bibitem{DKVanishing}
A.~D'Agnolo and M.~Kashiwara, \emph{Enhanced nearby and vanishing cycles in dimension one and
Fourier transform}, Publ. Res. Inst. Math. Sci. \textbf{59} (2023), 543--570.
\bibitem{Doucot2021}
J.~Dou\c{c}ot, \emph{Diagrams and irregular connections on the Riemann sphere},
arXiv:2107.02516v3.
\bibitem{Doucot2025}
J.~Dou\c{c}ot, \emph{New nonabelian Hodge graphs from twisted irregular connections},
arXiv:2509.24861.
\bibitem{Hiroe2017}
K.~Hiroe, \emph{Ramified irregular singularities of meromorphic connections and plane curve
singularities}, Josai Math. Monogr. \textbf{10} (2017), 161--192.
\bibitem{Hiroe2019}
K.~Hiroe, \emph{Index of rigidity of differential equations and Euler characteristic of their
spectral curves}, J. Geom. Phys. \textbf{162} (2021), 104060.
\bibitem{Martinez1994}
S.~Mart\'{i}nez, G.~Michon and J.~San Mart\'{i}n, \emph{Inverse of strictly ultrametric matrices
are of Stieltjes type}, SIAM J. Matrix Anal. Appl. \textbf{15} (1994), 98--106.
\bibitem{McDonald1995}
J.~J.~McDonald, M.~Neumann, H.~Schneider and M.~J.~Tsatsomeros, \emph{Inverse $M$-matrix
inequalities and generalized ultrametric matrices}, Linear Algebra Appl. \textbf{220} (1995),
321--341.
\bibitem{Mertens1874}
F.~Mertens, \emph{Ein Beitrag zur analytischen Zahlentheorie},
J. Reine Angew. Math. \textbf{78} (1874), 46--62.
\bibitem{Nabben1995}
R.~Nabben and R.~S.~Varga, \emph{Generalized ultrametric matrices---a class of inverse
$M$-matrices}, Linear Algebra Appl. \textbf{220} (1995), 365--390.
\bibitem{Sabbah}
C.~Sabbah, \emph{Introduction to Stokes Structures}, Lecture Notes in Mathematics 2060,
Springer, 2013.
\bibitem{Titchmarsh1986}
E.~C.~Titchmarsh, \emph{The Theory of the Riemann Zeta-Function}, 2nd ed., revised by
D.~R.~Heath-Brown, Oxford University Press, 1986.
\end{thebibliography}
\end{document}
